% !PRINTED 362
\weberpage{362}
%<<CONT
dern auch das Paar $x_3 y_3$ und folglich alle Paare $x_i y_i$ vorkommen, dass also dieser Complex von dem aus $x_1 y_1$ abgeleiteten überhaupt nicht verschieden ist. Wir haben also den Satz:

\gesperrt{2. Die Paare eines \gesperrt{Steiner}'schen Complexes haben die Eigenschaft, dass die acht Berührungspunkte von irgend zweien dieser Paare auf einem Kegelschnitte liegen, und dass man immer denselben Complex erhält, von welchem der sechs Paare man ausgehen mag.}

Hieraus ergiebt sich, dass drei Doppeltangenten eines \gesperrt{Steiner}'schen Complexes, wie $x_1, y_1, x_2$, von denen zwei ein Paar des Complexes bilden, ihre sechs Berührungspunkte auf einem Kegelschnitte haben. Dagegen giebt es wieder Systeme von drei Doppeltangenten (Tripel), deren sechs Berührungspunkte nicht auf einem Kegelschnitte liegen.

Nach einer von \gesperrt{Frobenius} eingeführten Bezeichnung bilden drei Doppeltangenten, wie $x_1, y_1, x_2$, deren Berührungspunkte auf einem Kegelschnitte liegen, ein \gesperrt{syzygetisches Tripel}. Drei Doppeltangenten, deren sechs Berührungspunkte nicht auf einem Kegelschnitte liegen, bilden ein \gesperrt{azygetisches Tripel}. Entsprechend wollen wir vier Doppeltangenten, deren acht Berührungspunkte auf einem Kegelschnitte liegen, ein \gesperrt{syzygetisches Quadrupel}, und irgend ein System von Doppeltangenten, von denen je drei azygetisch sind, ein \gesperrt{azygetisches System} nennen.

Hier gilt nun der folgende wichtige Satz:

\gesperrt{3. Drei Doppeltangenten, die in einem \gesperrt{Steiner}'schen Complexe vorkommen, so dass keine zwei von ihnen ein Paar bilden, wie $x_1, x_2, x_3$, sind immer azygetisch.}

Der Beweis dieses Satzes ergiebt sich einfach aus der Gleichungsform (9). Aus ihr ersieht man zunächst, dass die drei Linien $x_1, x_2, x_3$ sich nicht in einem Punkte schneiden; denn ein solcher Schnittpunkt würde auf der Curve $f$ liegen und wäre daher ein singulärer Punkt. Wir können also $x_1, x_2, x_3$ als Coordinaten einführen und demnach
\[
\begin{gathered}
y_1 = \alpha_1 x_1 + \alpha_2 x_2 + \alpha_3 x_3\\
y_2 = \beta_1 x_1 + \beta_2 x_2 + \beta_3 x_3\\
y_3 = \gamma_1 x_1 + \gamma_2 x_2 + \gamma_3 x_3
\end{gathered}
\]
%CONT>>
